\documentclass{article}
% NeurIPS 2025 format (same style as Holstege et al., NeurIPS 2025). Default option = anonymous submission.
% For an arXiv preprint use \usepackage[preprint]{neurips_2025}; camera-ready: [final].
\usepackage{neurips_2025}
\usepackage[utf8]{inputenc}
\usepackage[T1]{fontenc}
\usepackage{hyperref}
\usepackage{url}
\usepackage{booktabs}
\usepackage{amsfonts,amsmath,amssymb}
\usepackage{nicefrac}
\usepackage{microtype}
\usepackage{xcolor}
\usepackage{graphicx}
\usepackage{multirow}
\newif\ifdupfloats\dupfloatstrue

\title{Masking Cannot Remove What Is Inside the Region of Interest: Location-Dependent Shortcut Mitigation in Medical Imaging}

\author{Anonymous Author(s)\\Affiliation\\Address\\\texttt{email}}

\begin{document}
\maketitle

\begin{abstract}
Medical image classifiers exploit acquisition artifacts---hair, sonographer calipers, intestinal debris---that
correlate with the diagnosis; the standard defence restricts the model to the region of interest (ROI). We show that
this defence fails for artifacts inside the ROI, give a closed-form account of when and why, and evaluate
alternatives. In pre-registered, leakage-controlled traps built from real artifacts in four diseases, masking always
helps more for out-of-ROI than for in-ROI artifacts (crossover $+0.15$ to $+0.37$ reversed-test AUROC, all
Holm-significant), with frozen and fine-tuned backbones and in the zero-shot scores of vision--language foundation
models. Masking becomes harmful when it also removes disease context: in dermoscopy, controlled overlap sweeps,
fine-tuning (capsule $-0.28$) and unaltered test data (official thyroid split $-0.10$ on shortcut-conflicting cases;
two of three held-out dermoscopy hospitals). A linear-Gaussian model fitted only to each classifier's clean and
correlated performance matches its held-out reversed-test AUROC (mean absolute error 0.039) and the location crossover
across 14 cohort--backbone pairs ($r=0.997$). For frozen foundation-model features we propose mask-then-erase with
generic synthetic overlays (U-MtE), which needs no artifact examples, artifact masks or artifact labels---only the ROI
mask. It removes in-ROI shortcuts carried by distinct overlays (thyroid $+0.13$ to $+0.22$ over masking on three
backbones); erasure beats augmentation with the same overlays (thyroid, dermoscopy); a disease-protected variant prevents failure on
pathology-like artifacts; text-prompted erasure is a near no-op. Shortcuts carried by correlates of the artifact
(hair, chest drains; chest radiography is a boundary case) require label-based reweighting. Code, pre-registrations
and a backbone-agnostic tool are public.
\end{abstract}

\section{Introduction}\label{m:intro}
\begin{figure}[t]\centering\includegraphics[width=\linewidth]{figures/examples.pdf}
\caption{Real artifacts inside and outside the region of interest (green: ROI; red: artifact mask). Masking the
background removes an out-of-ROI artifact (2nd, 4th) but keeps an in-ROI one (1st, 3rd) while deleting the context
around it. Left: ovarian tumours (MMOTU); right: capsule endoscopy (SEE-AI).}\label{m:fig:teaser}\label{fig:examples}\end{figure}
Deep classifiers for medical images rely on features that correlate with the label but are not evidence of disease
\citep{geirhos2020shortcut,degrave2021ai}: surgical markings and rulers change melanoma predictions
\citep{winkler2019markings,winkler2021scalebars}, and pneumothorax classifiers detect the chest drain that treats the
condition \citep{oakden2020hidden,jimenez2023shortcuts}. The most common defence is to restrict the classifier to the
region of interest (ROI) by masking the background. Its rationale assumes that shortcuts live outside the ROI. Hair,
calipers and debris do not respect that assumption: when the artifact overlaps the ROI, masking cannot remove it---it
removes everything \emph{else}.

We ask whether the value of a mitigation depends on where the artifact sits, which quantity governs that dependence,
and how in-ROI shortcuts can be removed without localising the artifact. Our contributions:
\begin{enumerate}
\item \textbf{A location law for ROI masking}, established with pre-registered real-artifact traps in four diseases
(dermoscopic hair; thyroid and ovarian sonographer calipers; capsule-endoscopy debris), replicated with three frozen
backbones, two fine-tuned architectures, the zero-shot scores of vision--language models and on unaltered test data.
\item \textbf{A closed-form account} (Sec.~\ref{m:theory}) of when in-ROI masking harms, when erasure fails, and why
group reweighting is location-invariant, which matches held-out results of real encoders in a pre-registered test.
\item \textbf{U-MtE}, a representation-level remedy for frozen foundation-model features that needs no artifact
annotation (only the ROI mask), with evidence that erasure (not augmentation) is what works, a
disease-protected variant, and a delineation of the regimes in which it does not apply.
\item \textbf{Rigour}: pre-registration before every experiment, Holm-corrected primary claims (12/16 supported),
replication of an earlier pilot (73/88 numbers), and gaming-aware reporting.
\end{enumerate}

\section{Related work}\label{m:related}
\citet{bissoto2020debiasing} showed that artifact debiasing in dermoscopy is harder than it appears and built trap sets;
test-time selection \citep{bissoto2023tts} and generative repair \citep{maskmedpaint} address specific artifacts. Hair occlusion degrades HAM10000 classifiers \citep{wang2024hair}; adding synthetic dark-corner artifacts to training
improves classification \citep{pewton2024dca}; calipers are shortcuts even for fetal-ultrasound segmentation
\citep{lin2024segshortcut}; and classifiers on frozen foundation-model features show shortcut-like behaviour that
pooled evaluation hides \citep{germani2026aggregation}. None treats the artifact's location relative to the ROI as
a variable. Group-robust training
\citep{sagawa2020groupdro,kirichenko2023dfr,liu2021jtt} needs group labels or error sets; linear concept erasure
\citep{ravfogel2020inlp,belrose2023leace} and its task-preserving form SPLINCE \citep{splice2025} need concept labels;
biased-prompt projection \citep{chuang2023debiasing} derives directions from text. Random erasing and occlusion
augmentation \citep{zhong2020random,fong2019occlusions} use overlays as training noise. We estimate an artifact subspace
from synthetic \emph{insertion} pairs and remove it only after masking.

\section{Setup}\label{m:setup}
\paragraph{Traps} Let $Y$ be the diagnosis and $A$ the presence of an artifact. A trap trains with
$P(A{=}1\mid Y{=}1)=0.9$, $P(A{=}1\mid Y{=}0)=0.1$ and tests on a correlated (same rates), a reversed (swapped) and a
clean environment (artifact independent of the label: present in half of each class in real traps, absent in
sweeps). The overlap $r=|\text{artifact}\cap\text{ROI}|/|\text{artifact}|$ defines
Trap~A (in-ROI, $r\ge0.5$) and Trap~B (out-of-ROI, $r<0.1$); both share one artifact-free group, cells are matched on
source and label, and splits are grouped by lesion/patient identifiers where published, otherwise by perceptual hashes and frame
blocks (five seeds $\times$ five folds; a stricter embedding-based grouping is a sensitivity analysis). The primary metric is reversed-test AUROC; we always report correlated and clean AUROC, and summarise robustness
by $\min(\text{reversed},\text{correlated})$, which penalises methods that merely flip the shortcut. Uncertainty is a
hierarchical paired bootstrap (10{,}000 replicates, clusters = seeds or folds).
\paragraph{Cohorts} (Table~\ref{m:tab:data}; details in Appendix~\ref{sec:data}.) Heads are linear on frozen
DINOv2 ViT-B/14, MedSigLIP, DermLIP, RAD-DINO or ConvNeXt features; regularisation and threshold are chosen on clean
validation data only. Controlled sweeps paste a synthetic artifact (ruler, caliper, debris, tube) at overlap
$r\in\{0,0.25,0.5,0.75,1\}$.
\begin{table}[t]\centering\small
\caption{Real-artifact cohorts. Artifact masks: published (Wegley et al.), audited detector (95\% precision on 40
images), or expert-trained probe (pixel accuracy 0.92).}\label{m:tab:data}
\begin{tabular}{llrll}\toprule
Cohort & Modality & Images & Artifact (mask source) & ROI \\\midrule
ISIC 2019 & dermoscopy & 25{,}331 & hair/ruler (published masks) & lesion \\
TN3K/TNCD & thyroid ultrasound & 3{,}493 & calipers (rule-based detector) & nodule \\
MMOTU & ovarian ultrasound & 1{,}202 & calipers/overlays (same detector) & tumour \\
SEE-AI & capsule endoscopy & 5{,}481 & debris/bubbles (patch-token probe) & lesion boxes \\
NIH ChestX-ray14 & radiography & 29{,}687 & chest drains (image-level) & lungs \\\bottomrule
\end{tabular}\end{table}

\section{Mask-then-erase with a generic insertion library}\label{m:method}\label{sec:i2e}
Removing an in-ROI artifact requires its pixels; inserting one does not. A procedural generator $\mathcal{G}$
(strokes, bands, patches, blobs, dashed lines, glyphs; random colour, opacity and position, half centred in the ROI)
turns every masked training image $m(x_i)$ into a counterfactual $m(\mathcal{G}(x_i))$; the feature differences
$\Delta_i=f(m(\mathcal{G}(x_i)))-f(m(x_i))$ are independent of $Y$. U-MtE erases their span,
\begin{equation}
U_k=\mathrm{top}\text{-}k\ \mathrm{eigvecs}\big(\tfrac1n\textstyle\sum_i\Delta_i\Delta_i^\top\big),\qquad P(z)=z-U_kU_k^\top(z-\mu),
\end{equation}
with $k$ chosen without labels (90\% held-out difference energy), and fits the head on $P(f(m(x)))$. The generator
contains no ruler, hair, caliper or debris template. \emph{Disease protection} erases
$U_k'=\mathrm{orth}((I-WW^\top)U_k)$, where $W$ spans logistic label directions fitted on artifact-free training
images, guaranteeing $w^\top P(z)=w^\top z$ for $w\in\mathrm{span}(W)$. With image-level artifact labels the head is
group-balanced (U-MtE$_{\text{bal}}$).

\section{Why location matters}\label{m:theory}
Let the disease block have class-separation SNR $S$ and the artifact block SNR $A$ per unit of presence. For the
LDA-optimal linear head trained with $(p_1,p_0)$, the artifact's effective SNR is
$\tilde A=(p_1-p_0)^2A/(1+p_1(1-p_1)A)$ and
\begin{equation}
\mathrm{AUROC}_{\mathrm{rev}}=\Phi\!\left(\frac{S-\tilde A}{\sqrt{2(S+\tilde A)}}\right),
\end{equation}
increasing in $S$ and decreasing in $\tilde A$ (simulation agreement $\le0.009$ AUROC; Appendix~\ref{sec:theory}).
Masking maps $(S,\tilde A)\mapsto(S_m,\tilde A_m)$ with $\tilde A_m=0$ out of the ROI and $\tilde A_m\ge\tilde A$
inside. Hence in-ROI masking harms exactly when it removes disease context ($S_m<S$); the crossover is positive for
every $S_m$; erasure sets $\tilde A\to0$ but scales $S$ by $1-\rho^2$ ($\rho$: overlap of the erased subspace with the
disease direction), so it fails for pathology-like artifacts unless protected; and group reweighting yields
$\Phi(\sqrt{S/2})$ regardless of location (simulation in Appendix Fig.~\ref{fig:theory}).

\paragraph{Test on real encoders} The same head gives the AUROC of any test composition (Appendix~\ref{sec:theory}).
In a pre-registered test we fitted $(S,A)$ for every trained model to its clean and correlated AUROC only, with no
shared or tuned parameter, and compared the implied reversed-test AUROC with the held-out observation
(Fig.~\ref{m:fig:theory_predict}): mean absolute error 0.039 over 148 ERM and masking models ($r=0.92$; 0.108 for a
linear-symmetry heuristic). The implied location crossover agrees with the observed one across 14 cohort--backbone
pairs ($r=0.997$, error 0.014, 13/14 signs; the miss is a chest-radiograph cell at $-0.002$) and also for
fine-tuned networks (4/4 signs). One sign miss means that, by our pre-registered criterion, the model remains a
qualitative account; it fails only for chest drains, whose shortcut is carried by correlates that the clean and
correlated environments do not expose.
\begin{figure}[t]\centering\includegraphics[width=\linewidth]{figures/theory_predict.pdf}
\caption{Pre-registered test of the closed form on real encoders: $(S,A)$ fitted to each model's clean and correlated
AUROC, reversed-test results held out. (a) Reversed-test AUROC (ERM, masking). (b) Location crossover per cohort and
backbone (red: end-to-end fine-tuned, outside the model's assumptions). (c) Mask gain along controlled sweeps.}
\label{m:fig:theory_predict}\label{fig:theory_predict}\end{figure}

\section{Results}\label{m:results}
\subsection{A location law for ROI masking}
\begin{figure}[t]\centering\includegraphics[width=\linewidth]{figures/dose_response.pdf}
\caption{Controlled overlap sweeps: $\Delta$ reversed-test AUROC versus ERM as a synthetic artifact moves into the ROI
(95\% CIs). Masking (blue) declines with overlap except in chest radiography, where the encoder ignores extra-pulmonary
tubes; U-MtE with group balancing (green) is flat.}\label{m:fig:dose}\label{fig:dose}\end{figure}
\begin{table}[t]\centering\small
\caption{Masking gain (mask $-$ ERM, reversed AUROC) is always larger for out-of-ROI than in-ROI artifacts.
Crossover $=$ gain(out) $-$ gain(in), 95\% CI. Zero-shot: no task training. Natural: shortcut-conflicting pairs of
untouched test data.}\label{m:tab:law}
\begin{tabular}{llc}\toprule
Setting & Cohort (backbone) & Crossover [95\% CI] \\\midrule
\multirow{4}{*}{Real traps, frozen} & ISIC hair (DINOv2) & $+0.145$ [$+0.105$, $+0.181$] \\
 & Thyroid (DINOv2 / MedSigLIP / ConvNeXt) & $+0.230$ / $+0.163$ / $+0.284$ \\
 & Ovary (DINOv2) & $+0.170$ [$+0.116$, $+0.224$] \\
 & Capsule (DINOv2 / MedSigLIP / ConvNeXt) & $+0.368$ / $+0.332$ / $+0.292$ \\\midrule
\multirow{2}{*}{Fine-tuned} & Capsule (ResNet-50) & $+0.584$ [$+0.469$, $+0.691$] \\
 & Thyroid (ViT-S) & $+0.168$ [$+0.074$, $+0.265$] \\\midrule
\multirow{2}{*}{Zero-shot VLM} & Thyroid / Capsule (MedSigLIP) & $+0.103$ / $+0.049$ \\
 & Ovary (MedSigLIP) / Hair (DermLIP) & $+0.181$ / $+0.051$ \\\midrule
\multirow{2}{*}{Natural, mask $-$ ERM} & Thyroid (official split) & $-0.102$ [$-0.139$, $-0.066$] \\
 & ISIC held-out BCN / MSK & $-0.016$ / $-0.056$ \\\bottomrule
\end{tabular}\end{table}
In all four real-artifact cohorts the crossover is positive and survives Holm correction (Table~\ref{m:tab:law}); it
replicates with MedSigLIP and ConvNeXt, with end-to-end fine-tuning (ResNet-50 capsule; ViT-S thyroid), and in the
zero-shot scores of MedSigLIP and DermLIP, i.e.\ before any task training. Masking becomes \emph{harmful} for in-ROI
artifacts when it also removes disease context: real dermoscopic hair ($-0.037$ [$-0.057, -0.019$]), controlled sweeps
at full overlap (ovary $-0.197$ [$-0.292, -0.100$]; capsule $-0.102$ [$-0.162, -0.040$]; Fig.~\ref{m:fig:dose}),
fine-tuned capsule ($-0.284$ [$-0.347, -0.226$]) and on unaltered data, where masking lowers AUROC on
shortcut-conflicting cases for thyroid and two of three held-out ISIC hospitals (Table~\ref{m:tab:natural}); at the
validation-fixed operating point it halves thyroid sensitivity (0.734 $\to$ 0.347; AUPRC, Brier score and calibration
in Appendix Table~\ref{tab:natural_clinical}). The claims are robust to the
artifact-label definition (large-marker and strict-location variants) to a visual audit of the masks and to a stricter,
embedding-based leakage grouping (11 of 12 primary estimates unchanged; Appendix~\ref{sec:robust}). Chest radiography is a boundary case: the encoder ignores extra-pulmonary tubes, so
masking has nothing to remove outside the lungs.

\paragraph{Mechanism.} Three explanations were tested separately in the controlled dermoscopy sweep
(Appendix~\ref{sec:derm}). \emph{Occlusion} is not the cause: with the artifact present but uncorrelated with the
label, masking helps at every overlap. \emph{Lost context} alone is not the cause: dilating the mask at 50\% overlap
increases harm as it restores artifact pixels, and at full overlap harm is flat although the artifact's share of the
visible pixels falls threefold. The governing quantity is artifact \emph{retention}: masking keeps an in-ROI artifact,
and the masked model's counterfactual sensitivity to it roughly doubles.

\begin{figure}[t]\centering\includegraphics[width=0.78\linewidth]{figures/forest.pdf}
\caption{The 16 pre-registered primary tests with Holm correction (blue: supported, 12/16; red: not supported).}
\label{m:fig:forest}\label{fig:forest}\end{figure}

\begin{table}[t]\centering\small
\caption{Unaltered data (no resampling): change in AUROC on shortcut-conflicting pairs. Thyroid: official TN3K split;
ISIC: each hospital held out in turn; capsule: held-out frame blocks. 95\% CIs.}\label{m:tab:natural}
\begin{tabular}{lccc}\toprule
Test set & mask $-$ ERM & U-MtE $-$ mask & U-MtE$_{\text{bal}}$ $-$ mask \\\midrule
Thyroid (official split) & $-0.102$ [$-0.139$, $-0.066$] & $+0.051$ [$+0.033$, $+0.072$] & $+0.088$ [$+0.071$, $+0.108$] \\
ISIC, BCN held out & $-0.016$ [$-0.027$, $-0.005$] & $+0.003$ [$+0.000$, $+0.006$] & $+0.034$ [$+0.029$, $+0.039$] \\
ISIC, MSK held out & $-0.056$ [$-0.084$, $-0.029$] & $-0.008$ [$-0.015$, $-0.002$] & $+0.041$ [$+0.029$, $+0.053$] \\
ISIC, HAM held out & $+0.053$ [$+0.037$, $+0.068$] & $+0.023$ [$+0.017$, $+0.032$] & $+0.065$ [$+0.056$, $+0.075$] \\
Capsule (held-out frames) & $+0.021$ [$+0.000$, $+0.044$] & $+0.000$ [$-0.010$, $+0.009$] & $+0.014$ [$+0.003$, $+0.025$] \\\bottomrule
\end{tabular}\end{table}

\subsection{Removing in-ROI shortcuts}
\input{tables/main_inroi_compact}
Without any artifact example, artifact mask or label, U-MtE improves on masking where the artifact is a distinct overlay:
thyroid $+0.222$ [$+0.197, +0.251$] (DINOv2), $+0.204$ [$+0.174, +0.237$] (MedSigLIP), $+0.128$ [$+0.112, +0.145$]
(ConvNeXt); ovary $+0.215$ [$+0.182, +0.246$] (MedSigLIP); controlled sweeps $+0.08$ to $+0.28$; end-to-end
(ResNet-50, thyroid) $+0.143$ [$+0.099, +0.188$]. The same overlays used as training augmentation gain at most
$+0.03$; erasing their subspace gains $+0.207$ [$+0.182, +0.235$] more on thyroid. On capsule debris, which resembles
erosion fibrin, unprotected erasure removes disease signal; disease protection removes the failure. With image-level
artifact labels, U-MtE$_{\text{bal}}$ is the most robust model on thyroid with all three backbones
(Table~\ref{m:tab:remedies}; full values in Table~\ref{tab:main}).
Where the shortcut is carried by correlates of the artifact (hair co-varies with body site, age and sex; a drain comes
with a treated lung) no pixel- or insertion-based method helps much and DFR \citep{kirichenko2023dfr} is best.

\paragraph{Baselines} SPLINCE \citep{splice2025}, fitted with artifact labels, removes much disease signal when artifact
and label are strongly associated (clean AUROC 0.56--0.78) and flips the shortcut on capsule; U-MtE$_{\text{bal}}$ is
more robust on thyroid, capsule and hair (0.788 vs 0.596; 0.857 vs 0.683; 0.704 vs 0.564). Text-prompted erasure
(biased prompts, \citealp{chuang2023debiasing}) is a near no-op ($\le+0.008$ over masking): the direction along which a
\emph{described} artifact moves the text embedding is not the one along which the real artifact moves the image
embedding. Oracle-mask LaMa inpainting \citep{suvorov2022lama} is strong when artifact masks exist; protected U-MtE
matches it on thyroid without masks. Label-free, U-MtE beats JTT \citep{liu2021jtt} on thyroid ($+0.198$) and ovary
($+0.137$). A validation-driven selector beats masking in all five cohorts but not DFR everywhere (Appendix~\ref{sec:solution}).

\section{Discussion and limitations}\label{m:discussion}\label{sec:limits}
\input{sections/discussion}
\paragraph{Broader impact} Shortcut audits that report only overall accuracy can certify masked models that still
exploit in-ROI artifacts, and a mitigation that flips a shortcut can look like a success on a reversed test. We
recommend reporting artifact location relative to the ROI and performance on shortcut-conflicting cases together
with clean and correlated performance.

\section{Conclusion}
ROI masking is not neutral: it removes shortcuts outside the ROI and leaves---or amplifies---those inside it. This
holds across diseases, modalities, training regimes and foundation models, and a two-parameter model fitted to
clean and correlated performance matches the held-out results.
Erasing a subspace learned from generic synthetic overlays removes in-ROI shortcuts carried by distinct artifacts
without any annotation of the artifact.

\bibliographystyle{abbrvnat}
\bibliography{refs}

\appendix
\dupfloatsfalse
\input{sections/data}
\input{sections/theory}
\input{sections/results_derm}
\input{sections/results_multimodal}
\input{sections/results_robustness}
\input{sections/results_solution}
\section{Statistical plan and pre-registration}\label{app:stats}
Every experiment was specified in a dated, publicly versioned document before any model was fitted on its data. The 16
primary tests (location crossover, protected U-MtE vs masking, erase vs augment, U-MtE$_{\text{bal}}$ vs balanced; four
cohorts each) are Holm-corrected; 12/16 are supported (Fig.~\ref{fig:forest}). Everything else is secondary. Datasets
considered and not used (CANDID-PTX: no access; breast ultrasound: no artifact-free group; colonoscopy: no
reflection-free images; ISIC ink, cardiomegaly lines: too few images; pathology pen marks: unreliable masks) and the
leakage analysis are documented in the repository.
\section{Compute}\label{app:compute}
All experiments ran on one RTX 3090 (24 GB), 8 CPU cores and 31 GB RAM; the full pipeline takes about two days,
dominated by feature extraction and fine-tuning.

\section*{NeurIPS Paper Checklist}




\begin{enumerate}

\item {\bf Claims}
    \item[] Question: Do the main claims made in the abstract and introduction accurately reflect the paper's contributions and scope?
    \item[] Answer: \answerYes{} % Replace by \answerYes{}, \answerNo{}, or \answerNA{}.
    \item[] Justification: The abstract and Sec.~\ref{m:intro} state the location law, the theory, U-MtE and its limits; each claim is backed in Sec.~\ref{m:results} and Appendices \ref{sec:derm}--\ref{sec:solution}, including the negative results.
    \item[] Guidelines:
    \begin{itemize}
        \item The answer NA means that the abstract and introduction do not include the claims made in the paper.
        \item The abstract and/or introduction should clearly state the claims made, including the contributions made in the paper and important assumptions and limitations. A No or NA answer to this question will not be perceived well by the reviewers. 
        \item The claims made should match theoretical and experimental results, and reflect how much the results can be expected to generalize to other settings. 
        \item It is fine to include aspirational goals as motivation as long as it is clear that these goals are not attained by the paper. 
    \end{itemize}

\item {\bf Limitations}
    \item[] Question: Does the paper discuss the limitations of the work performed by the authors?
    \item[] Answer: \answerYes{} % Replace by \answerYes{}, \answerNo{}, or \answerNA{}.
    \item[] Justification: Sec.~\ref{m:discussion} (Limitations) and the regime analysis in Sec.~\ref{m:results} state where masking does not harm, where U-MtE fails and which datasets were infeasible.
    \item[] Guidelines:
    \begin{itemize}
        \item The answer NA means that the paper has no limitation while the answer No means that the paper has limitations, but those are not discussed in the paper. 
        \item The authors are encouraged to create a separate "Limitations" section in their paper.
        \item The paper should point out any strong assumptions and how robust the results are to violations of these assumptions (e.g., independence assumptions, noiseless settings, model well-specification, asymptotic approximations only holding locally). The authors should reflect on how these assumptions might be violated in practice and what the implications would be.
        \item The authors should reflect on the scope of the claims made, e.g., if the approach was only tested on a few datasets or with a few runs. In general, empirical results often depend on implicit assumptions, which should be articulated.
        \item The authors should reflect on the factors that influence the performance of the approach. For example, a facial recognition algorithm may perform poorly when image resolution is low or images are taken in low lighting. Or a speech-to-text system might not be used reliably to provide closed captions for online lectures because it fails to handle technical jargon.
        \item The authors should discuss the computational efficiency of the proposed algorithms and how they scale with dataset size.
        \item If applicable, the authors should discuss possible limitations of their approach to address problems of privacy and fairness.
        \item While the authors might fear that complete honesty about limitations might be used by reviewers as grounds for rejection, a worse outcome might be that reviewers discover limitations that aren't acknowledged in the paper. The authors should use their best judgment and recognize that individual actions in favor of transparency play an important role in developing norms that preserve the integrity of the community. Reviewers will be specifically instructed to not penalize honesty concerning limitations.
    \end{itemize}

\item {\bf Theory assumptions and proofs}
    \item[] Question: For each theoretical result, does the paper provide the full set of assumptions and a complete (and correct) proof?
    \item[] Answer: \answerYes{} % Replace by \answerYes{}, \answerNo{}, or \answerNA{}.
    \item[] Justification: Sec.~\ref{m:theory} states the assumptions (Gaussian class-conditional features, LDA-optimal linear head, additive artifact block); the derivation and a simulation check are in Appendix~\ref{sec:theory}.
    \item[] Guidelines:
    \begin{itemize}
        \item The answer NA means that the paper does not include theoretical results. 
        \item All the theorems, formulas, and proofs in the paper should be numbered and cross-referenced.
        \item All assumptions should be clearly stated or referenced in the statement of any theorems.
        \item The proofs can either appear in the main paper or the supplemental material, but if they appear in the supplemental material, the authors are encouraged to provide a short proof sketch to provide intuition. 
        \item Inversely, any informal proof provided in the core of the paper should be complemented by formal proofs provided in appendix or supplemental material.
        \item Theorems and Lemmas that the proof relies upon should be properly referenced. 
    \end{itemize}

    \item {\bf Experimental result reproducibility}
    \item[] Question: Does the paper fully disclose all the information needed to reproduce the main experimental results of the paper to the extent that it affects the main claims and/or conclusions of the paper (regardless of whether the code and data are provided or not)?
    \item[] Answer: \answerYes{} % Replace by \answerYes{}, \answerNo{}, or \answerNA{}.
    \item[] Justification: Protocol, splits, seeds, thresholds and arms are specified in Sec.~\ref{m:setup} and the appendices; every experiment has a dated pre-registration and a make target in the released code.
    \item[] Guidelines:
    \begin{itemize}
        \item The answer NA means that the paper does not include experiments.
        \item If the paper includes experiments, a No answer to this question will not be perceived well by the reviewers: Making the paper reproducible is important, regardless of whether the code and data are provided or not.
        \item If the contribution is a dataset and/or model, the authors should describe the steps taken to make their results reproducible or verifiable. 
        \item Depending on the contribution, reproducibility can be accomplished in various ways. For example, if the contribution is a novel architecture, describing the architecture fully might suffice, or if the contribution is a specific model and empirical evaluation, it may be necessary to either make it possible for others to replicate the model with the same dataset, or provide access to the model. In general. releasing code and data is often one good way to accomplish this, but reproducibility can also be provided via detailed instructions for how to replicate the results, access to a hosted model (e.g., in the case of a large language model), releasing of a model checkpoint, or other means that are appropriate to the research performed.
        \item While NeurIPS does not require releasing code, the conference does require all submissions to provide some reasonable avenue for reproducibility, which may depend on the nature of the contribution. For example
        \begin{enumerate}
            \item If the contribution is primarily a new algorithm, the paper should make it clear how to reproduce that algorithm.
            \item If the contribution is primarily a new model architecture, the paper should describe the architecture clearly and fully.
            \item If the contribution is a new model (e.g., a large language model), then there should either be a way to access this model for reproducing the results or a way to reproduce the model (e.g., with an open-source dataset or instructions for how to construct the dataset).
            \item We recognize that reproducibility may be tricky in some cases, in which case authors are welcome to describe the particular way they provide for reproducibility. In the case of closed-source models, it may be that access to the model is limited in some way (e.g., to registered users), but it should be possible for other researchers to have some path to reproducing or verifying the results.
        \end{enumerate}
    \end{itemize}


\item {\bf Open access to data and code}
    \item[] Question: Does the paper provide open access to the data and code, with sufficient instructions to faithfully reproduce the main experimental results, as described in supplemental material?
    \item[] Answer: \answerYes{} % Replace by \answerYes{}, \answerNo{}, or \answerNA{}.
    \item[] Justification: All datasets are public; code, frozen splits, pre-registrations and a two-command tool are released (anonymised link at submission).
    \item[] Guidelines:
    \begin{itemize}
        \item The answer NA means that paper does not include experiments requiring code.
        \item Please see the NeurIPS code and data submission guidelines (\url{https://nips.cc/public/guides/CodeSubmissionPolicy}) for more details.
        \item While we encourage the release of code and data, we understand that this might not be possible, so “No” is an acceptable answer. Papers cannot be rejected simply for not including code, unless this is central to the contribution (e.g., for a new open-source benchmark).
        \item The instructions should contain the exact command and environment needed to run to reproduce the results. See the NeurIPS code and data submission guidelines (\url{https://nips.cc/public/guides/CodeSubmissionPolicy}) for more details.
        \item The authors should provide instructions on data access and preparation, including how to access the raw data, preprocessed data, intermediate data, and generated data, etc.
        \item The authors should provide scripts to reproduce all experimental results for the new proposed method and baselines. If only a subset of experiments are reproducible, they should state which ones are omitted from the script and why.
        \item At submission time, to preserve anonymity, the authors should release anonymized versions (if applicable).
        \item Providing as much information as possible in supplemental material (appended to the paper) is recommended, but including URLs to data and code is permitted.
    \end{itemize}


\item {\bf Experimental setting/details}
    \item[] Question: Does the paper specify all the training and test details (e.g., data splits, hyperparameters, how they were chosen, type of optimizer, etc.) necessary to understand the results?
    \item[] Answer: \answerYes{} % Replace by \answerYes{}, \answerNo{}, or \answerNA{}.
    \item[] Justification: Sec.~\ref{m:setup} and Appendix~\ref{sec:data}: cohorts, backbones, regularisation grid, threshold selection on clean validation data, fine-tuning hyper-parameters.
    \item[] Guidelines:
    \begin{itemize}
        \item The answer NA means that the paper does not include experiments.
        \item The experimental setting should be presented in the core of the paper to a level of detail that is necessary to appreciate the results and make sense of them.
        \item The full details can be provided either with the code, in appendix, or as supplemental material.
    \end{itemize}

\item {\bf Experiment statistical significance}
    \item[] Question: Does the paper report error bars suitably and correctly defined or other appropriate information about the statistical significance of the experiments?
    \item[] Answer: \answerYes{} % Replace by \answerYes{}, \answerNo{}, or \answerNA{}.
    \item[] Justification: All effects are reported with 95\% hierarchical paired bootstrap CIs (10,000 replicates; clusters = seeds/folds); the primary family is Holm-corrected (Appendix~\ref{app:stats}).
    \item[] Guidelines:
    \begin{itemize}
        \item The answer NA means that the paper does not include experiments.
        \item The authors should answer "Yes" if the results are accompanied by error bars, confidence intervals, or statistical significance tests, at least for the experiments that support the main claims of the paper.
        \item The factors of variability that the error bars are capturing should be clearly stated (for example, train/test split, initialization, random drawing of some parameter, or overall run with given experimental conditions).
        \item The method for calculating the error bars should be explained (closed form formula, call to a library function, bootstrap, etc.)
        \item The assumptions made should be given (e.g., Normally distributed errors).
        \item It should be clear whether the error bar is the standard deviation or the standard error of the mean.
        \item It is OK to report 1-sigma error bars, but one should state it. The authors should preferably report a 2-sigma error bar than state that they have a 96\% CI, if the hypothesis of Normality of errors is not verified.
        \item For asymmetric distributions, the authors should be careful not to show in tables or figures symmetric error bars that would yield results that are out of range (e.g. negative error rates).
        \item If error bars are reported in tables or plots, The authors should explain in the text how they were calculated and reference the corresponding figures or tables in the text.
    \end{itemize}

\item {\bf Experiments compute resources}
    \item[] Question: For each experiment, does the paper provide sufficient information on the computer resources (type of compute workers, memory, time of execution) needed to reproduce the experiments?
    \item[] Answer: \answerYes{} % Replace by \answerYes{}, \answerNo{}, or \answerNA{}.
    \item[] Justification: Appendix~\ref{app:compute}: one RTX 3090, 8 CPU cores, 31 GB RAM; about two days end to end.
    \item[] Guidelines:
    \begin{itemize}
        \item The answer NA means that the paper does not include experiments.
        \item The paper should indicate the type of compute workers CPU or GPU, internal cluster, or cloud provider, including relevant memory and storage.
        \item The paper should provide the amount of compute required for each of the individual experimental runs as well as estimate the total compute. 
        \item The paper should disclose whether the full research project required more compute than the experiments reported in the paper (e.g., preliminary or failed experiments that didn't make it into the paper). 
    \end{itemize}
    
\item {\bf Code of ethics}
    \item[] Question: Does the research conducted in the paper conform, in every respect, with the NeurIPS Code of Ethics \url{https://neurips.cc/public/EthicsGuidelines}?
    \item[] Answer: \answerYes{} % Replace by \answerYes{}, \answerNo{}, or \answerNA{}.
    \item[] Justification: Only public, de-identified datasets under their licences; no human subjects; results reported whichever way they went.
    \item[] Guidelines:
    \begin{itemize}
        \item The answer NA means that the authors have not reviewed the NeurIPS Code of Ethics.
        \item If the authors answer No, they should explain the special circumstances that require a deviation from the Code of Ethics.
        \item The authors should make sure to preserve anonymity (e.g., if there is a special consideration due to laws or regulations in their jurisdiction).
    \end{itemize}


\item {\bf Broader impacts}
    \item[] Question: Does the paper discuss both potential positive societal impacts and negative societal impacts of the work performed?
    \item[] Answer: \answerYes{} % Replace by \answerYes{}, \answerNo{}, or \answerNA{}.
    \item[] Justification: Sec.~\ref{m:discussion}: masked models can be wrongly certified as shortcut-free; we recommend reporting artifact location and shortcut-conflicting performance. No foreseeable misuse beyond general model deployment risks.
    \item[] Guidelines:
    \begin{itemize}
        \item The answer NA means that there is no societal impact of the work performed.
        \item If the authors answer NA or No, they should explain why their work has no societal impact or why the paper does not address societal impact.
        \item Examples of negative societal impacts include potential malicious or unintended uses (e.g., disinformation, generating fake profiles, surveillance), fairness considerations (e.g., deployment of technologies that could make decisions that unfairly impact specific groups), privacy considerations, and security considerations.
        \item The conference expects that many papers will be foundational research and not tied to particular applications, let alone deployments. However, if there is a direct path to any negative applications, the authors should point it out. For example, it is legitimate to point out that an improvement in the quality of generative models could be used to generate deepfakes for disinformation. On the other hand, it is not needed to point out that a generic algorithm for optimizing neural networks could enable people to train models that generate Deepfakes faster.
        \item The authors should consider possible harms that could arise when the technology is being used as intended and functioning correctly, harms that could arise when the technology is being used as intended but gives incorrect results, and harms following from (intentional or unintentional) misuse of the technology.
        \item If there are negative societal impacts, the authors could also discuss possible mitigation strategies (e.g., gated release of models, providing defenses in addition to attacks, mechanisms for monitoring misuse, mechanisms to monitor how a system learns from feedback over time, improving the efficiency and accessibility of ML).
    \end{itemize}
    
\item {\bf Safeguards}
    \item[] Question: Does the paper describe safeguards that have been put in place for responsible release of data or models that have a high risk for misuse (e.g., pretrained language models, image generators, or scraped datasets)?
    \item[] Answer: \answerNA{} % Replace by \answerYes{}, \answerNo{}, or \answerNA{}.
    \item[] Justification: No high-risk models or data are released; the released artifacts are code and feature-level tools.
    \item[] Guidelines:
    \begin{itemize}
        \item The answer NA means that the paper poses no such risks.
        \item Released models that have a high risk for misuse or dual-use should be released with necessary safeguards to allow for controlled use of the model, for example by requiring that users adhere to usage guidelines or restrictions to access the model or implementing safety filters. 
        \item Datasets that have been scraped from the Internet could pose safety risks. The authors should describe how they avoided releasing unsafe images.
        \item We recognize that providing effective safeguards is challenging, and many papers do not require this, but we encourage authors to take this into account and make a best faith effort.
    \end{itemize}

\item {\bf Licenses for existing assets}
    \item[] Question: Are the creators or original owners of assets (e.g., code, data, models), used in the paper, properly credited and are the license and terms of use explicitly mentioned and properly respected?
    \item[] Answer: \answerYes{} % Replace by \answerYes{}, \answerNo{}, or \answerNA{}.
    \item[] Justification: Datasets are cited with licences in Appendix~\ref{sec:data} and the repository (ISIC CC BY-NC, MMOTU and SEE-AI CC BY 4.0, NIH ChestX-ray14, TN3K/TNCD, NEATX, RANZCR-CLiP); backbones are used under their model licences.
    \item[] Guidelines:
    \begin{itemize}
        \item The answer NA means that the paper does not use existing assets.
        \item The authors should cite the original paper that produced the code package or dataset.
        \item The authors should state which version of the asset is used and, if possible, include a URL.
        \item The name of the license (e.g., CC-BY 4.0) should be included for each asset.
        \item For scraped data from a particular source (e.g., website), the copyright and terms of service of that source should be provided.
        \item If assets are released, the license, copyright information, and terms of use in the package should be provided. For popular datasets, \url{paperswithcode.com/datasets} has curated licenses for some datasets. Their licensing guide can help determine the license of a dataset.
        \item For existing datasets that are re-packaged, both the original license and the license of the derived asset (if it has changed) should be provided.
        \item If this information is not available online, the authors are encouraged to reach out to the asset's creators.
    \end{itemize}

\item {\bf New assets}
    \item[] Question: Are new assets introduced in the paper well documented and is the documentation provided alongside the assets?
    \item[] Answer: \answerYes{} % Replace by \answerYes{}, \answerNo{}, or \answerNA{}.
    \item[] Justification: The code base and the U-MtE/SLAS tools are documented in the repository README with usage examples and tests.
    \item[] Guidelines:
    \begin{itemize}
        \item The answer NA means that the paper does not release new assets.
        \item Researchers should communicate the details of the dataset/code/model as part of their submissions via structured templates. This includes details about training, license, limitations, etc. 
        \item The paper should discuss whether and how consent was obtained from people whose asset is used.
        \item At submission time, remember to anonymize your assets (if applicable). You can either create an anonymized URL or include an anonymized zip file.
    \end{itemize}

\item {\bf Crowdsourcing and research with human subjects}
    \item[] Question: For crowdsourcing experiments and research with human subjects, does the paper include the full text of instructions given to participants and screenshots, if applicable, as well as details about compensation (if any)? 
    \item[] Answer: \answerNA{} % Replace by \answerYes{}, \answerNo{}, or \answerNA{}.
    \item[] Justification: No crowdsourcing or new human-subject data.
    \item[] Guidelines:
    \begin{itemize}
        \item The answer NA means that the paper does not involve crowdsourcing nor research with human subjects.
        \item Including this information in the supplemental material is fine, but if the main contribution of the paper involves human subjects, then as much detail as possible should be included in the main paper. 
        \item According to the NeurIPS Code of Ethics, workers involved in data collection, curation, or other labor should be paid at least the minimum wage in the country of the data collector. 
    \end{itemize}

\item {\bf Institutional review board (IRB) approvals or equivalent for research with human subjects}
    \item[] Question: Does the paper describe potential risks incurred by study participants, whether such risks were disclosed to the subjects, and whether Institutional Review Board (IRB) approvals (or an equivalent approval/review based on the requirements of your country or institution) were obtained?
    \item[] Answer: \answerNA{} % Replace by \answerYes{}, \answerNo{}, or \answerNA{}.
    \item[] Justification: Only publicly released, de-identified retrospective datasets are used; no new data were collected.
    \item[] Guidelines:
    \begin{itemize}
        \item The answer NA means that the paper does not involve crowdsourcing nor research with human subjects.
        \item Depending on the country in which research is conducted, IRB approval (or equivalent) may be required for any human subjects research. If you obtained IRB approval, you should clearly state this in the paper. 
        \item We recognize that the procedures for this may vary significantly between institutions and locations, and we expect authors to adhere to the NeurIPS Code of Ethics and the guidelines for their institution. 
        \item For initial submissions, do not include any information that would break anonymity (if applicable), such as the institution conducting the review.
    \end{itemize}

\item {\bf Declaration of LLM usage}
    \item[] Question: Does the paper describe the usage of LLMs if it is an important, original, or non-standard component of the core methods in this research? Note that if the LLM is used only for writing, editing, or formatting purposes and does not impact the core methodology, scientific rigorousness, or originality of the research, declaration is not required.
    %this research? 
    \item[] Answer: \answerNA{} % Replace by \answerYes{}, \answerNo{}, or \answerNA{}.
    \item[] Justification: LLMs are not a component of the proposed methods. (An AI coding assistant supported implementation and drafting under author supervision; authors should state this per venue policy.)
    \item[] Guidelines:
    \begin{itemize}
        \item The answer NA means that the core method development in this research does not involve LLMs as any important, original, or non-standard components.
        \item Please refer to our LLM policy (\url{https://neurips.cc/Conferences/2025/LLM}) for what should or should not be described.
    \end{itemize}

\end{enumerate}



\end{document}
